% CCECE 2027 - "Grounding the AI Scrum Master", IEEE conference format (6 pages max incl. references).
% Companion to the sprint-risk paper; separate question, separate results.
% Pilot code, prompts, raw agent answers and results: pilot/ in this folder.
% Rev. 8 Nov 2026: kappa, past-lateness rule and checkpoint stability from pilot/results/rule_baseline.json;
% closing-on-date finding reframed against the Scrum Guide; four reference entries corrected; AI-use note added.
% Rev. 8 Nov 2026 (evening): synced with the companion sprint-risk paper's audit. Shuffle claim replaced by the
% leakage case studies; model AUC given within projects too; Lee et al. author list completed; short-lag and weekday evidence for the label; spillover advice now names the count baseline.
\documentclass[conference]{IEEEtran}
\usepackage[T1]{fontenc}
\usepackage[utf8]{inputenc}
\usepackage{amsmath,amssymb}
\usepackage{array}
\usepackage{url}
\usepackage{microtype}
\begin{document}
\title{Grounding the AI Scrum Master:\\ A Sprint-Replay Test and Pilot}
\author{\IEEEauthorblockN{Hani AbuSharkh}
\IEEEauthorblockA{Higher Colleges of Technology\\Al Ain, UAE\\habusharkh@hct.ac.ae}}
\maketitle

\begin{abstract}
Agent systems that carry the Scrum Master label have started to appear, yet none we found has been judged against what actually happened in real sprints, while work that predicts sprint outcomes from Jira histories stops at prediction and never acts. We make four contributions. First, we map each Scrum Master accountability to the best published evidence and show that the label has run ahead of it. Second, we name two failure modes that bound the design space: an ungrounded agent says things the data does not support, and an over-grounded one repeats its model and adds nothing. Third, we propose sprint replay, a hindsight-free test on public data: rebuild the board as it stood, run a grounded agent against a model-only arm and an agent-only arm, and report agreement and a cost-weighted measure of what grounding adds. Fourth, we run a pilot on 40 held-out TAWOS sprints at two checkpoints. A Claude agent given the risk model agreed with it on 91\% of decisions (Cohen's kappa 0.82), with no detectable cost difference at any price of a warning we tried. The same agent without the model made the same call as a one-line rule on past lateness for 39 of 40 sprints. One warning in five advised closing the sprint on its planned date and carrying unfinished work over. Scrum asks for exactly that, which shows that late closure measures timebox discipline more than delivery. The test is cheap, it already catches the over-grounded failure, and it exposed a weak target label.
\end{abstract}

\begin{IEEEkeywords}
agile software development, evaluation, grounding, LLM agents, Scrum, sprint analytics
\end{IEEEkeywords}

\section{Introduction}

Much of what a Scrum Master does never reaches the tracker. Shastri \emph{et al.}~\cite{shastri} name six everyday functions of the role: facilitating, mentoring, negotiating, adapting processes, coordinating and protecting the team. Few of them leave much of a trace in Jira; most happen in conversation.

Language-model agents have nonetheless started to appear in Scrum roles. Since 2024, published systems have given an agent the Scrum Master role in a code-generation pipeline~\cite{agilecoder}, simulated agile teams with role agents~\cite{cognisim}, and turned burndown charts into status reports under the name of an AI Scrum Master~\cite{alliata}.

Our position is that the label has run ahead of the evidence, and that a specific, affordable test can bring it back into line. The test is sprint replay: rebuild a project's board as it stood at a chosen moment, give the agent only what was knowable then, let it decide whether and what to say, and score that against what the sprint went on to do. No partner organisation or live team is needed. What is needed is a risk model the agent can cite, a simple rule to check the model against, and the discipline to report how often the agent merely repeats either of them.

This paper makes four contributions:
\begin{itemize}
\item a map of the role's accountabilities against the best published evidence, which shows the label running ahead of it (Section~\ref{sec:gap});
\item two documented failure modes, ungrounded and over-grounded, that bound what a credible system must avoid (Section~\ref{sec:fail});
\item sprint replay, a hindsight-free test with a model-only arm, an agent-only arm, an agreement rate and a cost-weighted measure of added value (Sections~\ref{sec:ground} and~\ref{sec:replay});
\item a pilot of that test on 40 public sprints, which already shows one of the two failure modes and a flaw in the target label (Section~\ref{sec:pilot}).
\end{itemize}

\section{The Label and the Evidence}
\label{sec:gap}

Systems that use Scrum roles fall into three families. The first organises code generation: AgileCoder's Scrum Master agent reviews a backlog, picks sprint tasks and judges deliverability, and is scored on code benchmarks alone~\cite{agilecoder}. The role supplies a vocabulary, not a set of responsibilities. The second simulates agile teams; CogniSim models scaled-agile roles with language-model agents but has no Scrum Master agent at all, and its own future work asks for predictive analytics for delays~\cite{cognisim}. The third targets the role directly, and here the evidence is real but small. An action-research study put assistants into daily Scrums and refinement for three teams at a large organisation: about a third of early refinement outputs were wrong because the input data was stale, and participants asked that it inform rather than warn~\cite{cabrero}. Another study prompts models to write status reports from burndown charts and scores the text with similarity metrics, with no agent and no prediction~\cite{alliata}.

Software management accounts for under 1\% of language-model usages in a review of 395 software-engineering papers~\cite{hou}, and a 2026 review of AI in Scrum finds no applications in the implementation or review-and-retrospect phases~\cite{peras}. Table~\ref{tab:gap} sets the role's accountabilities against the best evidence we found. It is a structured reading of these reviews and the studies they point to, not a systematic review of its own.

\begin{table}[t]
\centering\footnotesize
\caption{Scrum Master accountabilities against the best evidence we found. Rows without a citation are absences: we found no study, and the reviews in the text report none.}
\label{tab:gap}
\setlength{\tabcolsep}{3pt}
\begin{tabular}{>{\raggedright\arraybackslash}p{0.36\linewidth} >{\raggedright\arraybackslash}p{0.17\linewidth} >{\raggedright\arraybackslash}p{0.37\linewidth}}
\hline
Accountability & Coverage & Best evidence\\
\hline
Effort and story-point estimation & Well covered & Real data, many projects~\cite{deepse}\\
Delivery-risk prediction at planning & Moderate & Real Jira histories~\cite{tse2018,sprint2vec}\\
In-sprint monitoring & Thin & Checkpoint prediction, no agent~\cite{tse2018}\\
Daily Scrum facilitation & Thin & One field site, three teams~\cite{cabrero}\\
Impediment detection and removal & Near empty & No removal study found\\
Review and stakeholder reporting & Near empty & Text-similarity scoring~\cite{alliata}\\
Retrospectives, coaching, team health & Near empty & None in a 2026 review~\cite{peras}\\
End-to-end agent in the role & Label only & Code benchmarks and demos~\cite{agilecoder}\\
\hline
\end{tabular}
\end{table}

Three patterns follow. The more relational the task, the thinner the evidence: of the six functions Shastri \emph{et al.}~\cite{shastri} identify, only facilitation and coordination have any. Evaluation habits are weak, with small samples, single runs and author-labelled ground truth, and where run-to-run stability has been measured it is poor~\cite{yoon}. And the data is lopsided. Public sprint histories are plentiful: the TAWOS dataset of agile open-source projects alone holds thousands of sprints and millions of change-log entries~\cite{tawos}, while public transcripts of stand-ups, reviews and retrospectives are, as far as we can tell, absent. The half of the role a tracker records is therefore the half in which an agent can be held to account today.

\section{Two Ways to Fail}
\label{sec:fail}

\paragraph{Ungrounded} Asked certification-style Scrum questions, GPT-5 recalled well and interpreted less well; at least one prompt got 16\% of the items wrong, with the weaker results on interpretive topics~\cite{psm}. Generated sprint plans were judged unfit for direct use, with priority accuracy swinging from 44\% to 92\% across runs~\cite{yoon}. In mission-critical risk analysis, retrieval-grounded models had the lowest hallucination rates~\cite{esem24}. In this role the characteristic failure is a confident wrong statement, and the characteristic fix is grounding.

\paragraph{Over-grounded} The opposite failure is less discussed and, for research, more dangerous. Given a predictive model as a tool, agents tend to copy it: in one study they reproduced a graph model's output on 97.6\% to 99.2\% of cases, and stronger backbones deferred more, not less~\cite{tool}. An agent that repeats its model is the model with added latency. Any claim of an agent in this role should therefore report its agreement with the raw model and include a model-only arm.

Neither failure is particular to Scrum. The same two questions arise wherever an agent sits on top of a predictive model: when should it believe the model, and when is acting on it worth the cost? Sprints are an unusually good place to study both, because outcomes are public, time-stamped, and free of hindsight once the board is rebuilt.

\section{What Grounding Means Here}
\label{sec:ground}

Grounding means that the agent's statements about risk trace back to models fitted on real sprint histories, under an evaluation that does not leak the future. Some sprint-outcome studies report accuracies near or above 0.9 using features visible only once the sprint is under way~\cite{perez}, or after resampling with no stated split order~\cite{obike}, and leakage taxonomies now list software engineering among the affected fields~\cite{leakage}. A stricter pipeline produces lower numbers, and the lower numbers are the trustworthy ones. Both mistakes are easy to reproduce. In our own study on public Jira data~\cite{artifact}, adding a sprint's own completion ratio, known only once it ends, lifts the area under the ROC curve (AUC) for spillover from 0.83 to 1.00. Oversampling the whole pool before splitting lifts it for late closure, a sprint closed in Jira more than a day after its planned end, from 0.64 to 0.69. An agent that cites an inflated model inherits the inflation and adds a confident voice to it.

Grounding also limits what the agent may suggest. In defect prediction, guidance derived from explanation methods often fails to fix what it predicts and differs from what developers actually change~\cite{saner}, while restricting suggestions to changes with historical precedent brings them much closer to real practice~\cite{counteract}. Carried into sprints, an agent should not propose cutting scope merely because the model's expectation falls; it should propose changes that comparable teams actually made.

What is left for the agent is narrower than the label and more interesting than the model: reading the issue text the tabular model never sees, deciding whether a situation is worth raising, choosing the moment, and wording it so that a team can act without feeling audited. Each of these is measurable.

\section{Sprint Replay as the Bar}
\label{sec:replay}

\paragraph{Construction} Public datasets record field-level changes with timestamps~\cite{tawos}, which is enough to rebuild a board as it stood at any point in a sprint. For sprint $s$ and checkpoint $c \in \{0, 0.25, 0.5, 0.75\}$ of the planned duration, let $B_{s,c}$ be the board rebuilt from events time-stamped no later than that point. Every arm sees $B_{s,c}$ and nothing later.

\paragraph{Arms} Three at minimum: the risk model alone ($M$), the language-model agent without the model ($L$), and the grounded agent ($G$). Without $M$, no one can tell whether the agent added anything; without $L$, no one can tell whether grounding did. A rule a Scrum Master could apply by hand, such as the team's recent lateness, is a useful floor for both. At each replay point, arm $a$ emits a decision $d_a(s,c) \in \{\textit{warn}, \textit{silent}\}$ and, for $L$ and $G$, the message it would send.

\paragraph{Agreement} Over the set $P$ of replay points, for each agent arm $a \in \{L, G\}$,
\[
A_a = \frac{1}{|P|} \sum_{(s,c) \in P} \mathbf{1}\big[d_a(s,c) = d_M(s,c)\big],
\]
reported as a headline result together with Cohen's kappa, which removes the agreement two arms would reach by chance at their warning rates. $A_G$ is the one that matters; $A_L$ shows how much of it the board alone explains.

\paragraph{Added value} We borrow the cost structure of alarm-based prescriptive process monitoring~\cite{alarm}. A warning costs $\kappa$, a failed sprint costs $F$, and a warning raised at checkpoint $c$ averts a fraction $e(c)$ of that cost, with $e$ falling as the sprint proceeds. Let $y_s = 1$ if sprint $s$ failed, $w_a(s) = 1$ if arm $a$ warned at any checkpoint, and $c^{*}_a(s)$ the first checkpoint at which it did. Then
\begin{gather*}
C_a(s) = \kappa\, w_a(s) + F\, y_s \big(1 - e(c^{*}_a(s))\, w_a(s)\big),\\
\Delta = \min\big(\bar{C}_M, \bar{C}_L\big) - \bar{C}_G ,
\end{gather*}
where $\bar{C}_a$ is the mean of $C_a(s)$ over sprints. A positive $\Delta$ means the grounded agent beats the better of the two baselines. Because $\kappa/F$ and $e$ differ between teams, $\Delta$ is reported as a curve over the cost ratio rather than as one number. Read together, a high $A_G$ with $\Delta \approx 0$ says the agent is the model, while $\Delta > 0$ with $A_G$ well below one is evidence of judgement.

\paragraph{Reliability and calibration} Language-model arms are rerun and their variance reported, and any probability an arm states is checked for calibration. So are alarm thresholds: in our study above~\cite{artifact}, one threshold caught 11\% to 86\% of late sprints across eight outside projects.

\paragraph{What replay cannot do} Replay measures what an agent would have said, not whether saying it would have changed the sprint. The one controlled study of language-model facilitation found it drew out more information, mostly from quieter participants, but not better decisions~\cite{alsobay}. Replay is a filter, not a verdict: a system that fails it should not reach a team, and one that passes has earned a study, not a deployment.

\section{A Pilot Replay}
\label{sec:pilot}

To show that the test is cheap, we ran it at small scale on late closure.

\paragraph{Setup} We drew 40 sprints from the held-out part of TAWOS~\cite{tawos} (the last 30\% of each project's history): 20 that closed late and 20 that did not, among the 472 held-out sprints that had committed work and were still open at mid-sprint (seed 42). Five projects supply 37 of the 40, and one supplies 11. Each sprint was replayed at planning ($c=0$) and halfway ($c=0.5$), so every arm made 80 decisions. Arm $M$ is gradient boosting from our leakage-aware pipeline~\cite{artifact} on the same data, trained only on the first 70\% of each project: planning features at $c=0$ (held-out AUC 0.751 on all 786 test sprints, but 0.634 when scored within each project) and planning plus progress features at $c=0.5$ (0.757 on 775). Much of its skill therefore tells late projects from punctual ones. It warns when its score passes the threshold that maximises sensitivity plus specificity on out-of-fold training scores. Arms $L$ and $G$ are Claude (model claude-opus-5-5), each given the rebuilt board as text: team history, the committed issues with their type and story points and, at $c=0.5$, which issues were resolved, moved out or added. Project names, issue keys and calendar dates are hidden. $G$ also sees $M$'s probability, threshold and call, and is told it may overrule the model. Boards went out in two batches of 20 per arm and checkpoint, in random order. The agents ran as sub-agents in the Claude app, not through the API, so sampling settings were not under our control. Each agent's only tool call was the read of its own prompt file; the run logs show no other call. The eight agent runs took about 20 seconds each and about 0.55 million tokens in all, most of it reading the boards. Each answer gives a warn or silent decision, a probability, and the message the agent would send. Costs use $e(0)=0.5$ and $e(0.5)=0.25$ and are reweighted to the 28\% late rate of the 472 eligible sprints. Confidence intervals (CI) come from 2,000 bootstrap resamples of the 40 sprints.

\begin{table}[t]
\centering\footnotesize
\caption{Pilot replay on 40 held-out TAWOS sprints (20 late), per arm and checkpoint. $A$: share of decisions equal to the model's. $F_1$ is at the sample's 50\% late rate and would be lower at the true 28\%.}
\label{tab:pilot}
\setlength{\tabcolsep}{4pt}
\begin{tabular}{llcccccc}
\hline
$c$ & Arm & Warn rate & Sens. & Spec. & $F_1$ & AUC & $A$\\
\hline
0 & $M$ model & .58 & .75 & .60 & .70 & .72 & 1.00\\
0 & $L$ agent & .40 & .55 & .75 & .61 & .69 & .73\\
0 & $G$ grounded & .60 & .75 & .55 & .68 & .74 & .93\\
0.5 & $M$ model & .65 & .75 & .45 & .65 & .71 & 1.00\\
0.5 & $L$ agent & .43 & .60 & .75 & .65 & .70 & .73\\
0.5 & $G$ grounded & .60 & .75 & .55 & .68 & .72 & .90\\
\hline
\end{tabular}
\end{table}

\paragraph{Agreement and added value} Table~\ref{tab:pilot} gives the result. The grounded agent agreed with the model on 91\% of replay points (95\% CI 84--98\%), a Cohen's kappa of 0.82. It overruled the model 7 times in 80 and was right in 4 of them, too few to tell judgement from noise. The agent without the model agreed with it on 73\% (60--85\%; kappa 0.47), and 19 of its 22 departures were silences where the model warned. That arm turns out to be a simple rule. Warning whenever at least half of the team's last three sprints closed late reproduces its planning-time call on 39 of 40 sprints. The rule ranks sprints with an AUC of 0.63, and the model's edge over it, 0.09, has an interval from $-0.08$ to $0.25$. Over all 786 held-out sprints, the recent late rate behind this rule ranks a project's own sprints at chance (0.51 within projects), so the rule mostly recognises late projects. Calls also barely moved between checkpoints: $M$ changed its call on 3 of the 40 sprints, $G$ on 2 and $L$ on 1, so the 80 decisions per arm carry little more information than 40. On ranking, no arm separates from another: every AUC difference from $M$ is within 0.03, with intervals that reach $-$0.20 on the low side. On cost, we cannot tell $G$ and $M$ apart at any price of a warning from $\kappa/F = 0.02$ to $0.5$: the gap never exceeds $0.02F$ per sprint, and every interval contains zero. The more cautious $L$ arm is cheaper than $G$ once a warning costs a fifth of a late close or more, so $\Delta$ is negative from $\kappa/F = 0.2$ upward, and clearly so at $0.5$, where it is $-0.086F$ $[-0.161, -0.019]$. In the terms of Section~\ref{sec:replay}, this is a high $A_G$ with $\Delta$ near zero or below. On these boards, the grounded agent added nothing the model did not already say, and the ungrounded agent added nothing past lateness did not already say.

\paragraph{What the agents advised} The agent arms sent 81 warnings. One author read all of them and coded each for one thing: whether it advised closing the sprint on its planned date whatever was left unfinished (the coded file is in the artifact; a second coder would be needed to report agreement). Most warnings asked the team to trim scope or re-plan, citing the board's own numbers. But 16 (8 in each arm) advised closing on the date and carrying unfinished work over. In Scrum terms that advice is sound: a Sprint is a fixed timebox, and unfinished work goes back to the Product Backlog~\cite{scrumguide}. It is also the cheapest way to avoid the late-closure label, since it moves the closing date and not the work done. The wider data agree. Across all 2,602 sprints, 40\% of late closures are one to three days late, and sprints planned to end on a Friday or a weekend close late 57\% of the time, against 28\% otherwise~\cite{artifact}. The agents found that the label rewards timebox discipline, which a Scrum Master should want, but says little about delivery, which is what a risk warning is for. A replay that scored decisions alone would have missed this. Late closure is a weak sole target for an agent in this role. It should be scored together with the share of committed work left unfinished. A yes-or-no spillover label will not do, because counting the committed issues alone ranks it at 0.88~\cite{artifact}.

\paragraph{Limits of the pilot} This is one run of one model on 40 sprints, with a balanced sample, wide intervals and most sprints from five projects. It falls short of the bar in Section~\ref{sec:replay} on two counts: one run instead of several, and no calibration check on the stated probabilities. We did not vary the prompt either, so part of $A_G$ may be the prompt's framing of the model's call. We did not check the numbers the agents quoted against the boards, so the pilot says nothing yet about the ungrounded failure. The TAWOS extract we used has no issue titles or descriptions, so the agents could not use the text that a tabular model never sees, which is where we expect an agent to add most. Issue type is the final recorded value, story points are as of the checkpoint, and $e(c)$ is assumed. The pilot shows that the test runs, cheaply, and that it can catch an agent repeating its model or a simple rule. It does not show that agents cannot add value.

\section{Limits}

Public sprint data comes from open-source projects, and enterprise teams may behave differently. Models may have seen public tracker data in pre-training, which is one more reason to run the $L$ arm, since an agent that recalls outcomes would score well without any model, and to replicate outside software. The relational half of the role stays out of reach while no public ceremony corpus exists. Finally, autonomy: developers keep planning and management work at low autonomy levels~\cite{choudhuri}, and the one field study in this role reports participants asking to be informed rather than warned~\cite{cabrero}. An agent built to inform is easier to accept, and to evaluate honestly, than one built to decide.

\section{Conclusion}

The AI Scrum Master exists as a label and a few pilots; no system yet has had its judgement measured against real sprints, while sprint-prediction work senses without acting. Joining the two can be tested cheaply: replay real sprints, ground the agent in models whose evaluation does not leak the future, and report how often the agent repeated its model, or a simple rule, and what it added when it did not. Our pilot did this on 40 sprints at two checkpoints. The grounded agent agreed with its model nine times in ten, we could not tell their costs apart, and the agent without the model behaved like a rule on past lateness. Its advice also showed that late closure can be met by closing on time, as Scrum asks anyway. The next replay should score the share of work delivered as well as closing dates, check every arm against a count of open issues, and add issue text, more runs and more sprints. That is where an agent would have to earn its place.

% TODO (Hani): confirm this wording is accurate before submission. IEEE asks authors to disclose
% AI-generated content (text, code, figures) in the acknowledgments.
\section*{Acknowledgment}
Claude (Anthropic) served as the agent under test in the pilot and was also used to help write the pilot code and edit the text. The author checked the reported numbers against the result files and is responsible for the content.

\begin{thebibliography}{25}
\bibitem{shastri} Y. Shastri, R. Hoda, and R. Amor, ``Spearheading agile: The role of the scrum master in agile projects,'' \emph{Empir. Softw. Eng.}, vol. 26, no. 1, p. 3, 2021.
\bibitem{agilecoder} M. H. Nguyen, T. P. Chau, P. X. Nguyen, and N. D. Q. Bui, ``AgileCoder: Dynamic collaborative agents for software development based on agile methodology,'' in \emph{Proc. IEEE/ACM Int. Conf. AI Foundation Models and Software Engineering (FORGE)}, 2025.
\bibitem{cognisim} K. Cinkusz and J. A. Chudziak, ``Agile software management with cognitive multi-agent systems,'' in \emph{Proc. 17th Int. Conf. Agents and Artificial Intelligence (ICAART)}, 2025.
\bibitem{alliata} Z. Alliata, T. Singhal, and A.-M. Bozagiu, ``The AI scrum master: Using large language models (LLMs) to automate agile project management tasks,'' in \emph{Agile Processes in Software Engineering and Extreme Programming -- Workshops (XP 2024)}, LNBIP, vol. 524. Cham: Springer, 2025, pp. 110--122.
\bibitem{cabrero} B. Cabrero-Daniel \emph{et al.}, ``Exploring human-AI collaboration in agile: Customised LLM meeting assistants,'' in \emph{Agile Processes in Software Engineering and Extreme Programming (XP 2024)}, LNBIP. Springer, 2024.
\bibitem{hou} X. Hou \emph{et al.}, ``Large language models for software engineering: A systematic literature review,'' \emph{ACM Trans. Softw. Eng. Methodol.}, vol. 33, no. 8, 2024.
\bibitem{peras} D. Peras, L. de-Marcos, and Z. Stapi\'c, ``AI-enhanced Scrum: A literature review on the impact of artificial intelligence on Scrum processes,'' \emph{Informatica}, vol. 50, no. 12, 2026.
\bibitem{deepse} M. Choetkiertikul, H. K. Dam, T. Tran, T. Pham, A. Ghose, and T. Menzies, ``A deep learning model for estimating story points,'' \emph{IEEE Trans. Softw. Eng.}, vol. 45, no. 7, pp. 637--656, 2019.
\bibitem{tse2018} M. Choetkiertikul, H. K. Dam, T. Tran, A. Ghose, and J. Grundy, ``Predicting delivery capability in iterative software development,'' \emph{IEEE Trans. Softw. Eng.}, vol. 44, no. 6, pp. 551--573, 2018.
\bibitem{sprint2vec} M. Choetkiertikul, P. Banyongrakkul, C. Ragkhitwetsagul, S. Tuarob, H. K. Dam, and T. Sunetnanta, ``Sprint2Vec: A deep characterization of sprints in iterative software development,'' \emph{IEEE Trans. Softw. Eng.}, vol. 51, no. 1, pp. 220--242, 2025.
\bibitem{yoon} Y. Yoon, K. Iwan, M. Zwart, X. Qin, H. Lee, and M. Spichkova, ``Scrum sprint planning: LLM-based and algorithmic solutions,'' 2025, arXiv:2512.18966.
\bibitem{tawos} V. Tawosi, A. Al-Subaihin, R. Moussa, and F. Sarro, ``A versatile dataset of agile open source software projects,'' in \emph{Proc. 19th Int. Conf. Mining Software Repositories (MSR)}, 2022, pp. 707--711.
\bibitem{psm} M. Perkusich \emph{et al.}, ``Prompting GPT-5 on Scrum certification questions: An empirical accuracy study,'' 2026, arXiv:2607.00049.
\bibitem{esem24} M. Esposito, F. Palagiano, V. Lenarduzzi, and D. Taibi, ``Beyond words: On large language models actionability in mission-critical risk analysis,'' in \emph{Proc. 18th ACM/IEEE Int. Symp. Empirical Software Engineering and Measurement (ESEM)}, 2024.
\bibitem{tool} Z. Wang and P. Vemuri, ``When the tool decides: LLM agents defer blindly to graph neural network tools, and stronger backbones defer more,'' 2026, arXiv:2606.14476.
\bibitem{perez} Y. J. P\'erez Castillo, S. D. Orantes Jim\'enez, and P. O. Letelier Torres, ``Sprint management in Agile approach: Progress and velocity evaluation applying machine learning,'' \emph{Information}, vol. 15, no. 11, p. 726, 2024.
\bibitem{obike} P. G. Obike, V. E. Ekong, and O. U. Obot, ``A novel scoring-based Agile sprint deliverability prediction and prioritization using machine learning,'' \emph{Int. J. Comput. Appl.}, vol. 187, no. 49, pp. 40--53, 2025.
\bibitem{leakage} S. Kapoor and A. Narayanan, ``Leakage and the reproducibility crisis in machine-learning-based science,'' \emph{Patterns}, vol. 4, no. 9, p. 100804, 2023.
\bibitem{artifact} H. AbuSharkh, ``Sprint risk prediction from reconstructed Jira histories: Code, data and results,'' GitHub repository, 2026. [Online]. Available: \url{https://github.com/hani-ziad/SwEng_202610}
\bibitem{saner} G. Lee, H. Ju, and S. U.-J. Lee, ``Can we trust the actionable guidance from explainable AI techniques in defect prediction?'' in \emph{Proc. IEEE Int. Conf. Software Analysis, Evolution and Reengineering (SANER)}, 2025.
\bibitem{counteract} K. Oueslati, G. Laberge, M. Lamothe, and F. Khomh, ``Mining action rules for defect reduction planning,'' in \emph{Proc. ACM Int. Conf. Foundations of Software Engineering (FSE)}, 2024.
\bibitem{alarm} S. A. Fahrenkrog-Petersen, N. Tax, I. Teinemaa, M. Dumas, M. de Leoni, F. M. Maggi, and M. Weidlich, ``Fire now, fire later: Alarm-based systems for prescriptive process monitoring,'' \emph{Knowl. Inf. Syst.}, vol. 64, no. 2, pp. 559--587, 2022.
\bibitem{alsobay} M. Alsobay, D. M. Rothschild, J. M. Hofman, and D. G. Goldstein, ``Bringing everyone to the table: An experimental study of LLM-facilitated group decision making,'' 2025, arXiv:2508.08242.
\bibitem{scrumguide} K. Schwaber and J. Sutherland, ``The Scrum Guide,'' 2020. [Online]. Available: \url{https://scrumguides.org}
\bibitem{choudhuri} R. Choudhuri, C. Bird, C. Badea, M. Gerosa, and A. Sarma, ``You shall not pass! Where and why developers draw the line on AI autonomy,'' 2026, arXiv:2607.00533.
\end{thebibliography}

\end{document}
